\documentclass[12pt]{article}
\usepackage{amsmath,amssymb,amsfonts}
\usepackage[margin=1in]{geometry}
\begin{document}
\begin{center}
{\large\bfseries 31st Irish Mathematical Olympiad\\12 May 2018}
\end{center}
\bigskip\noindent\textbf{Paper 1}\bigskip
\begin{enumerate}
\item Mary and Pat play the following number game. Mary picks an initial integer greater than 2017. She then multiplies this number by 2017 and adds 2 to the result. Pat will add 2019 to this new number and it will again be Mary's turn. Both players will continue to take alternating turns. Mary will always multiply the current number by 2017 and add 2 to the result when it is her turn. Pat will always add 2019 to the current number when it is his turn. Pat wins if one of the numbers obtained is divisible by 2018. Mary wants to prevent Pat from winning the game. Determine, with proof, the smallest initial integer Mary could choose in order to achieve this.

\item The triangle $ABC$ is right-angled at $A$. Its incentre is $I$, and $H$ is the foot of the perpendicular from $I$ on $AB$. The perpendicular from $H$ on $BC$ meets $BC$ at $E$, and it meets the bisector of $\angle ABC$ at $D$. The perpendicular from $A$ on $BC$ meets $BC$ at $F$. Prove that $\angle EFD = 45^\circ$.

\item Find all functions $f(x) = ax^2 + bx + c$, with $a \neq 0$, such that
\[
f(f(1)) = f(f(0)) = f(f(-1)).
\]

\item We say that a rectangle with side lengths $a$ and $b$ \emph{fits inside} a rectangle with side lengths $c$ and $d$ if either ($a \leq c$ and $b \leq d$) or ($a \leq d$ and $b \leq c$). For instance, a rectangle with side lengths 1 and 5 fits inside a rectangle with side lengths 6 and 2. Suppose $S$ is a set of 2019 rectangles, all with integer side lengths between 1 and 2018 inclusive. Show that there are three rectangles $A$, $B$, and $C$ in $S$ such that $A$ fits inside $B$, and $B$ fits inside $C$.

\item Points $A$, $B$ and $P$ lie on the circumference of a circle $\Omega_1$ such that $\angle APB$ is an obtuse angle. Let $Q$ be the foot of the perpendicular from $P$ on $AB$. A second circle $\Omega_2$ is drawn with centre $P$ and radius $PQ$. The tangents from $A$ and $B$ to $\Omega_2$ intersect $\Omega_1$ at $F$ and $H$ respectively. Prove that $FH$ is tangent to $\Omega_2$.
\end{enumerate}

\newpage
\begin{center}
{\large\bfseries 31st Irish Mathematical Olympiad\\12 May 2018}
\end{center}
\bigskip\noindent\textbf{Paper 2}\bigskip
\begin{enumerate}
\setcounter{enumi}{5}
\item Find all real-valued functions $f$ satisfying
\[
f(2x + f(y)) + f(f(y)) \;=\; 4x + 8y
\]
for all real numbers $x$ and $y$.

\item Let $a$, $b$, $c$ be the side lengths of a triangle. Prove that
\[
2(a^3 + b^3 + c^3) \;<\; (a + b + c)(a^2 + b^2 + c^2) \;\leq\; 3(a^3 + b^3 + c^3).
\]

\item Let $M$ be the midpoint of side $BC$ of an equilateral triangle $ABC$. The point $D$ is on $CA$ extended such that $A$ is between $D$ and $C$. The point $E$ is on $AB$ extended such that $B$ is between $A$ and $E$, and $|MD| = |ME|$. The point $F$ is the intersection of $MD$ and $AB$. Prove that $\angle BFM = \angle BME$.

\item The sequence of positive integers $a_1, a_2, a_3, \ldots$ satisfies
\[
a_{n+1} = a_n^2 + 2018 \qquad \text{for } n \geq 1.
\]
Prove that there exists at most one $n$ for which $a_n$ is the cube of an integer.

\item The game of \emph{Greed} starts with an initial configuration of one or more piles of stones. Player 1 and Player 2 take turns to remove stones, beginning with Player 1. At each turn, a player has two choices:
\begin{itemize}
\item take one stone from any one of the piles (a simple move);
\item take one stone from each of the remaining piles (a greedy move).
\end{itemize}
The player who takes the last stone wins. Consider the following two initial configurations:
\begin{enumerate}
\item[(a)] There are 2018 piles, with either 20 or 18 stones in each pile.
\item[(b)] There are four piles, with 17, 18, 19, and 20 stones, respectively.
\end{enumerate}
In each case, find an appropriate strategy that guarantees victory to one of the players.
\end{enumerate}
\end{document}
